\section{Fundamentals}


%%% Fundamental Equation Y = C + I + G 

\begin{frame}{Stock and Flow Quantities}
    \mytext{
    Economics is concerned with economic activity. 
    
    We are concerned with \textbf{income:} 
    \begin{itemize}
        \item the production (of income),
        \item the use (of income) in
        \begin{itemize}
            \item consumption, or 
            \item savings and investment.
        \end{itemize}
    \end{itemize}
    
    Pure stocks are less a concern. Anything `sitting' around does not economically count. It is only interesting when sold or bought, when it yields some kind of interest, when it grows or shrinks.
    }
    \mycomment{
        Economic activity concerns the production, distribution or consumption of goods. Looking at a picture or a sunset and enjoying it is not economic activity. Living in a spacious house is not an economic activity. Living in a cave is not an economic activity.

        Stocks are interesting insofar they have value; but the value of economic concern is the exchange value, not the (individual) use value. For economics it does not matter how much you enjoy your car; economics asks only how much you can sell it for.
    }
\end{frame}

\begin{frame}{Consumption and Production}
    \mytext{
    Most of what we consume is produced immediately before it is consumed. 

    What is produced `appears'; what is consumed `disappears'.

    Products (and services) come into existence through the combination of \textbf{production factors.}
    }
    \mycomment{
    }
\end{frame}

\begin{frame}{Value Added}
    \showmytikz{../shared/tikz/tikz-valueadded}
    \mycomment{Production is more than just re-combination. The product is more than the sum of its parts.

    Typical production factors are:
    \begin{itemize}
        \item raw materials,
        \item pre-products or parts,
        \item energy,
        \item labour,
        \item the means of production: tools, machines, factory buildings etc.,
        \item know-how.
    \end{itemize}

    }
\end{frame}

\begin{frame}{Value Added}
    \showmytikz[2]{../shared/tikz/tikz-valueadded}
    \mycomment{

    This \textbf{value added} originates from the way the factors are combined; the final product has more value (on markets) than the costs (of the factors). 
    
    The factors used in production are sourced on markets; the product is sold on a market. All this makes economically only sense as long as value is added (or at least not destroyed).
    }
\end{frame}

\begin{frame}{Production Factors}
    \showmytikz[3]{../shared/tikz/tikz-valueadded}
    \mycomment{
        Some factors are simply material or energy inputs which are bought on markets. 

        These can be stored to some extent, but are usually processed immediately.       

        These factors are remunerated by payment.
    }
\end{frame}

\begin{frame}{Production Factors}
    \showmytikz[4]{../shared/tikz/tikz-valueadded}
    \mycomment{
        Labour is another important factor. No factory runs without someone pressing the buttons. Labour is very different from material inputs, as it is sourced in different ways. Typically work contracts are longer employment contracts. 

        Labour is remunerated by wages or salaries.
    }
\end{frame}

\begin{frame}{Production Factors}
    \showmytikz[5]{../shared/tikz/tikz-valueadded}
    \mycomment{
        Finally, some factors remain at the site, and can even be accumulated. These are the means of production: tools, machines, buildings, land and other properties. 

        Another word for these assets is \textbf{capital.} They are tangible assets (such as tools) and intangible assets (such as know-how, brands, or patents).

        Capital is remunerated by \textbf{rents or interest rates} paid to their owners or for the loans that were required to purchase them.
    }
\end{frame}



%%% Circular Model of the Economy

\begin{frame}{Circular Model of the Economy}
    \showmytikz[1]{../shared/tikz/tikz-flowmodel-own}
\end{frame}

\begin{frame}{Circular Model of the Economy}
    \showmytikz[1,2]{../shared/tikz/tikz-flowmodel-own}
\end{frame}

\begin{frame}{Circular Model of the Economy}
    \showmytikz[1,2,3]{../shared/tikz/tikz-flowmodel-own}
\end{frame}

\begin{frame}{Circular Model of the Economy}
    \showmytikz[1,2,3,4]{../shared/tikz/tikz-flowmodel-own}
\end{frame}


%%% Gross Domestic Product

\begin{frame}{Gross Domestic Product}
    \mytext{
    The sum of everything produced in an Economy is the Gross Domestic Product.

    It can be expressed or calculated in three ways:
    \begin{enumerate}
        \item by Origin
        \item by Distribution
        \item by Expenditure
    \end{enumerate}

    Note that it does not express assets or changes of ownership.
    }
\end{frame}


\begin{frame}{Gross Domestic Product by Origin}
    \mycomment{
    Registering Gross Value Added (GVA) in products and services across various sectors shows where GDP has been created. Net taxes have to be added, since this is value re-distributed from the private sector to the state.

    The distribution analysis helps understanding the structure of production in an economy. 
    }
    \showmytikz[1]{../shared/tikz/tikz-gdp}
    
\end{frame}

\begin{frame}{Gross Domestic Product by Distribution}
    \mycomment{
    Registering how the income generated is distributed along different groups and in different types of income delivers good insights on the distributive effects.
    }
    \showmytikz[2]{../shared/tikz/tikz-gdp}
    
\end{frame}

\begin{frame}{Gross Domestic Product by Expenditure}
    \mycomment{
    Everything produced (all supply) is also consumed, or saved (demand). For economic analysis, this expenditure statistics is the most important statistic.
    }
    \showmytikz[3]{../shared/tikz/tikz-gdp}
    
\end{frame}

\begin{frame}{Gross Domestic Product}
    \mycomment{
    There are three ways to calculate the Gross Domestic Product by summing up accounts kept by national statistic bureaus. Each of them must result in the same value, be it not for errors and omissions.

    The actual composition, and the development of this composition, of each of these balances is of course of major interest. In our figure, the entries are not proportional. 
    }
    \showmytikz[1,2,3]{../shared/tikz/tikz-gdp}
    
\end{frame}

\begin{frame}{GDP Germany}
    \showmybitmap{shared/bitmaps/BIP2025destatis.png}
\end{frame}


\begin{frame}{GDP by Origin}
    \mytext{
    The Equation of Origin shows Gross Value Added (GVA) $+$ taxes on products $-$ subsidies on products:
    \begin{equation}        
    \begin{split}
        \text{GDP} = & + \sum \left(\text{GVA}\right)\\ 
              & + \text{taxes on products}\\
              & - \text{subsidies on products}
    \end{split}
    \end{equation}
    with
    \begin{equation}
        \begin{split}
            \text{GVA} = & + \text{Gross Value of Output}\\
                         & - \text{Intermediate Consumption}
        \end{split}
    \end{equation}
    }
\end{frame}

\begin{frame}{GDP by Distribution}
    \mytext{
    Everything produced (GDP) is distributed as someone or somebodies' income; so as salaries or wages, as profits kept in the companies, or taxes minus subsidies.
    \begin{equation}
        \begin{split}
            \text{GDP} = & + \text{compensation of employees} \\
                         & + \text{gross operating surplus \& mixed income} \\
                         & + \text{taxes on products}\\
                         & - \text{subsidies on products}
        \end{split}        
    \end{equation}
    }
\end{frame}

\begin{frame}{GDP by Distribution}
    \mytext{
    This can be broken down further:
    \begin{equation}
        \begin{split}
            \text{GDP} = & + \text{wages} \\
                     = & + \text{rents} \\
                     = & + \text{interest} \\
                     = & + \text{profits} \\
                     = & + \text{taxes} \\
                     = & - \text{subsidies} \\
                     = & + \text{depreciation} \\
                     = & + \text{net foreign factor income}  
        \end{split}        
    \end{equation}
    }
\end{frame}


\begin{frame}{GDP by Expenditure}
    \mytext{
    The most interesting approach looks at how everything produced (GDP) is used. 
    \begin{equation}
        \begin{aligned}
            & Y & \text{GDP} = & + \text{consumption} && C\\
            & &              & + \text{investment} && I\\
            & &              & + \text{government spending} && G \\
            & &              & + \text{exports} && Ex \\
            & &              & - \text{imports} && Im 
        \end{aligned}
    \end{equation}
        
    }    
\end{frame}

\begin{frame}{GDP by Expenditure}
    \mytext{
    The Fundamental Expenditure Equation in Macroeconomics:
        \begin{equation}
            \boxed{Y = C + I + G + Ex - Im}
        \end{equation}

        We can also write:
        \begin{equation}
           Y = C + S + G + Ex - Im
        \end{equation}

        This is because of the fundamental identity of savings and investment:
        \begin{equation}
            S = I
        \end{equation}


    }    
\end{frame}



\begin{frame}{Production and Consumption}
    \mytext{
        Everything that is produced is either consumed or saved:

        \begin{equation}
            S = Y - C = I
        \end{equation}

        (abstracting, for now, from government expenditure and from net exports.)

        This is a fundamental balance / identity in respect to \textbf{net} savings and \textbf{net} investment, and relates to \textbf{real} assets.

        Transfer of property or the creation of debt is an entirely different issue. If Alice spends not (hence `saves') half of her income, but lends it to Bob instead (hence `gives Bob a credit' or `buys a bond issued by Bob'), then that income is transferred to Bob, who in turn can either consume it (as if it was a consumer credit) or invest it into real assets.
    }
\end{frame}
 

\begin{frame}{Keynesian Cross}
    \showmytikz{../shared/tikz/tikz-keynesiancross}
    \mycomment{
    Aggregate demand $AD$ is consumption, investment, and government expense;
    \begin{equation}
        AD = C+I+G
    \end{equation}
    with
    \begin{equation}
        C=a+b\left(Y-t\right)
    \end{equation}
    and $I,G,t$ constant; so that
    \begin{equation}
        AD = e + b\cdot Y
    \end{equation}
    with $G$ part of the autonomous part $e$.
    }
\end{frame}

\begin{frame}{Keynesian Cross}
    \showmytikz[2]{../shared/tikz/tikz-keynesiancross}
    \mycomment{
    An increase of government spending increases aggregate demand, shifting the $AD$ curve up. The result is a higher output $Y$.

    The same applies for autonomous increase of $I$.
    }
\end{frame}

\begin{frame}{Fiscal Multiplier}
    \mycomment{
    The idea with the multiplier is that any extra inserted expenditure (by fiscal policy, but anything else for that matter) is someone else's income, who will then re-spend a part of that, saving another part. 

    The size of the multiplier will depend on the propensity to spend (the inverse of the propensity to save), and on taxes paid each time.
    }
    \showmytikz{../shared/tikz/tikz-multiplier}
\end{frame}

\begin{frame}{Fiscal Multiplier}
    \mycomment{
    For now, let us abstract from taxes and net exports.
    }
    \mytext{
    Consumption $C$ is a (constant) part of income $Y$:
    \begin{equation}
        C=C_0+c\cdot Y
    \end{equation}
    Thus
    \begin{equation}
        Y=C_0+c\cdot Y + I + G
    \end{equation}
    For (in the short run) constant $I,C_0$ we get
    \begin{equation}
        \Delta Y = \frac{\Delta G}{1-c}=\frac{1}{s}
    \end{equation}
    with $s$ the savings quota. --
    Thus the multplier is
    \begin{equation}
        M=\frac{1}{1-c}=\frac{1}{s}.
    \end{equation}
    }
    \mycomment{
    % There are a number of (simplifying) assumptions here. One is that net exports $NX=Ex-Im$ be zero. Sin
    }
\end{frame}

\begin{frame}{Fiscal Multiplier: Critique}
    \mytext{
    \begin{itemize}
        \item crowding out private demand
        \item government debt means higher taxes:
        \item[$\rightarrow$] subjects reduce spending
        \item debt financing raises interest rates
        \item monetary policy counteracts
        \item lock-in problem
        \item might depend on government debt level
    \end{itemize}
    }
\end{frame}